\documentclass[11pt,a4paper]{article}

% ---------- Codificación e idioma ----------
\usepackage[utf8]{inputenc}
\usepackage[T1]{fontenc}
\usepackage[spanish,es-tabla,es-nodecimaldot]{babel}
\usepackage{lmodern}
\usepackage{amssymb}

% ---------- Geometría y layout ----------
\usepackage[a4paper,margin=2.5cm,headheight=22pt]{geometry}
\usepackage{parskip}
\usepackage{fancyhdr}
\usepackage{titlesec}
\usepackage{needspace}
\emergencystretch=3em
\sloppy

% ---------- Color y gráficos ----------
\usepackage{graphicx}
\usepackage{xcolor}
\usepackage[most]{tcolorbox}
\usepackage{float}
\usepackage{caption}

% ---------- Tablas ----------
\usepackage{tabularx}
\usepackage{longtable}
\usepackage{booktabs}
\usepackage{colortbl}
\usepackage{array}

% ---------- Listas, código y enlaces ----------
\usepackage{enumitem}
\usepackage{fancyvrb}
\usepackage{hyperref}

% ---------- Paleta de marca DermaUH ----------
\definecolor{dermaPrimary}{HTML}{A21CAF}
\definecolor{dermaPrimaryLight}{HTML}{D946EF}
\definecolor{dermaSoft}{HTML}{FAE8FF}
\definecolor{dermaBg}{HTML}{FDF4FF}
\definecolor{dermaText}{HTML}{4A044E}
\definecolor{dermaTextSoft}{HTML}{7E2F7D}
\definecolor{adminColor}{HTML}{B2416E}
\definecolor{noteColor}{HTML}{0F766E}
\definecolor{warnColor}{HTML}{B45309}
\definecolor{codeBg}{HTML}{F6F2F7}

\hypersetup{
    colorlinks=true,
    linkcolor=dermaPrimary,
    urlcolor=dermaPrimary,
    citecolor=dermaPrimary,
    pdftitle={Manual de Usuario - DermaUH Image Database},
    pdfauthor={Equipo DermaUH},
    pdfsubject={Manual de Usuario},
    bookmarksopen=true
}

% ---------- Estilo de títulos ----------
\titleformat{\section}
  {\color{dermaPrimary}\Large\bfseries}
  {\thesection}{0.6em}{}
\titleformat{\subsection}
  {\color{dermaTextSoft}\large\bfseries}
  {\thesubsection}{0.6em}{}
\titleformat{\subsubsection}
  {\color{dermaTextSoft}\normalsize\bfseries}
  {\thesubsubsection}{0.6em}{}

\titlespacing*{\section}{0pt}{1.4em}{0.8em}
\titlespacing*{\subsection}{0pt}{1.1em}{0.5em}

% Cada sección principal comienza en una página nueva.
\newcommand{\seccion}[1]{\clearpage\section{#1}}

% ---------- Encabezado / pie de página ----------
\pagestyle{fancy}
\fancyhf{}
\renewcommand{\headrulewidth}{0.4pt}
\renewcommand{\headrule}{\hbox to\headwidth{\color{dermaPrimaryLight}\leaders\hrule height \headrulewidth\hfill}}
\fancyhead[L]{\small\color{dermaTextSoft} Manual de Usuario}
\fancyhead[R]{\small\color{dermaTextSoft} DermaUH Image Database}
\fancyfoot[L]{\small\color{dermaTextSoft} Versión 1.0}
\fancyfoot[R]{\small\color{dermaTextSoft} Página \thepage}

\fancypagestyle{plain}{%
  \fancyhf{}
  \fancyfoot[R]{\small\color{dermaTextSoft} Página \thepage}
  \renewcommand{\headrulewidth}{0pt}
}

% ---------- Figuras ----------
\captionsetup{font=small, labelfont={bf,color=dermaPrimary}, justification=centering}
\captionsetup[figure]{name=Figura}

% \captura[ancho]{archivo}{pie de figura}: inserta una captura de pantalla
% desde figuras/, con marco fino y etiqueta fig:<archivo>.
\newcommand{\captura}[3][0.96\textwidth]{%
  \begin{figure}[H]
    \centering
    {\setlength{\fboxsep}{0pt}\setlength{\fboxrule}{0.5pt}%
     \fcolorbox{black!25}{white}{%
       \includegraphics[width=#1,height=0.68\textheight,keepaspectratio]{figuras/#2}}}
    \caption{#3}
    \label{fig:#2}
  \end{figure}}

% ---------- Cajas informativas ----------
\newtcolorbox{notabox}{
  colback=dermaSoft, colframe=dermaPrimaryLight, boxrule=0.6pt,
  arc=2pt, left=8pt, right=8pt, top=6pt, bottom=6pt,
  fonttitle=\bfseries, title=NOTA
}
\newtcolorbox{adminbox}{
  colback=white!95!adminColor, colframe=adminColor, boxrule=0.6pt,
  arc=2pt, left=8pt, right=8pt, top=6pt, bottom=6pt,
  fonttitle=\bfseries\color{white}, coltitle=white,
  colbacktitle=adminColor, title=SOLO ADMINISTRADORES
}
\newtcolorbox{warnbox}{
  colback=white!93!warnColor, colframe=warnColor, boxrule=0.6pt,
  arc=2pt, left=8pt, right=8pt, top=6pt, bottom=6pt,
  fonttitle=\bfseries, title=IMPORTANTE
}
\newtcolorbox{tipbox}{
  colback=white!93!noteColor, colframe=noteColor, boxrule=0.6pt,
  arc=2pt, left=8pt, right=8pt, top=6pt, bottom=6pt,
  fonttitle=\bfseries, title=CONSEJO
}

% Bloques de comandos de consola (anexo técnico).
\DefineVerbatimEnvironment{consola}{Verbatim}{%
  fontsize=\small, frame=single, framerule=0.4pt, rulecolor=\color{black!25},
  framesep=3mm, xleftmargin=2mm, xrightmargin=2mm}

% ---------- Convenciones tipográficas ----------
% Nombre de un campo de formulario
\newcommand{\campo}[1]{\textbf{\textcolor{dermaTextSoft}{#1}}}
% Botón, enlace o casilla de la interfaz
\newcommand{\boton}[1]{\colorbox{dermaSoft}{\textcolor{dermaPrimary}{\textbf{\small #1}}}}
% Rol de usuario
\newcommand{\rol}[1]{\textcolor{dermaPrimary}{\textbf{#1}}}
% Opción de menú o ruta de navegación
\newcommand{\menu}[1]{\textsf{\textbf{#1}}}
% Mensaje literal mostrado por el sistema
\newcommand{\mensaje}[1]{\textit{«#1»}}
% Separador de rutas de menú
\newcommand{\flecha}{\,$\triangleright$\,}

\newcolumntype{Y}{>{\raggedright\arraybackslash}X}

\setlist[itemize]{leftmargin=1.4em, itemsep=2pt, topsep=4pt}
\setlist[enumerate]{leftmargin=1.6em, itemsep=2pt, topsep=4pt}
% Lista para procedimientos paso a paso
\newlist{pasos}{enumerate}{1}
\setlist[pasos]{label=\textbf{\textcolor{dermaPrimary}{Paso \arabic*.}}, leftmargin=4.2em,
  labelwidth=3.8em, align=left, itemsep=3pt, topsep=4pt}

\begin{document}

% ==========================================================
% PORTADA
% ==========================================================
\begin{titlepage}
  \centering
  \vspace*{1.5cm}
  \includegraphics[width=5cm]{assets/logo-dermauh.png}\\[1.8cm]
  {\Huge\bfseries\color{dermaPrimary} Manual de Usuario\\[0.4cm]}
  {\LARGE\color{dermaTextSoft} DermaUH Image Database\\[0.3cm]}
  \vspace{0.4cm}
  {\large\color{dermaTextSoft} Plataforma para el almacenamiento, gestión y consulta\\
   de imágenes dermatológicas clínicas}
  \vfill
  {\color{dermaPrimaryLight}\rule{0.6\textwidth}{0.6pt}}\\[0.5cm]
  {\large Versión del documento: 1.0}\\
  {\large Septiembre de 2026}
  \vspace{1.5cm}
\end{titlepage}

% ==========================================================
% CONTROL DEL DOCUMENTO
% ==========================================================
\section*{Control del documento}

\begin{center}
\begin{tabularx}{\textwidth}{@{}l l Y@{}}
\toprule
\rowcolor{dermaSoft}
\textbf{Versión} & \textbf{Fecha} & \textbf{Descripción del cambio} \\
\midrule
1.0 & Septiembre de 2026 & Versión inicial del manual de usuario, con capturas de
pantalla de la plataforma y anexo de despliegue con Docker. \\
\bottomrule
\end{tabularx}
\end{center}

\vspace{1em}
\begin{notabox}
Las capturas de pantalla de este manual se obtuvieron de una instalación de
demostración con \textbf{datos ficticios}: las instituciones, personas y casos
clínicos que aparecen en ellas son ilustrativos y no corresponden a pacientes ni
a centros reales. Según la configuración y los datos de su instalación, las
pantallas pueden presentar diferencias menores.
\end{notabox}

\clearpage
\tableofcontents
\clearpage
\listoffigures

% ==========================================================
\seccion{Introducción}
% ==========================================================

\subsection{Propósito del manual}

Este documento explica cómo utilizar la aplicación web \textbf{DermaUH Image
Database}. Describe, paso a paso y con capturas de pantalla, cada una de las
tareas que los usuarios pueden realizar: crear una cuenta, iniciar sesión,
explorar y descargar imágenes, consultar los análisis estadísticos, gestionar el
perfil personal y, en el caso de los administradores, gestionar usuarios y
solicitudes.

\subsection{Alcance}

El manual cubre todas las funciones disponibles en la interfaz web de la
plataforma. La integración con otros sistemas mediante la API se presenta de forma
resumida en la sección~\ref{sec:desarrolladores}; su referencia técnica completa
está disponible dentro de la propia aplicación. La instalación de la plataforma se
describe en el Anexo~\ref{anexo:despliegue}, dirigido al personal técnico.

\subsection{Destinatarios}

\begin{itemize}
  \item \textbf{Especialistas clínicos e investigadores} que consultan el
  repositorio de imágenes, sus metadatos y sus estadísticas.
  \item \textbf{Administradores de la plataforma}, responsables de aprobar cuentas,
  autorizar descargas y gestionar usuarios.
  \item \textbf{Equipos técnicos} que integran otras aplicaciones con la plataforma
  o que se encargan de su despliegue.
\end{itemize}

No se requieren conocimientos informáticos especiales para las secciones
dirigidas a usuarios finales; basta con saber utilizar un navegador web.

\subsection{Cómo usar este manual}

\begin{itemize}
  \item Si es la primera vez que utiliza la plataforma, lea en orden las secciones
  \ref{sec:descripcion} a \ref{sec:primeros-pasos}.
  \item Para realizar una tarea concreta, localícela en el índice: cada tarea se
  describe como un \textbf{procedimiento numerado}.
  \item Si aparece un mensaje que no entiende o algo no funciona como espera,
  consulte la sección~\ref{sec:problemas}.
  \item Los términos médicos y técnicos se explican en el
  glosario (sección~\ref{sec:glosario}).
\end{itemize}

\subsection{Convenciones tipográficas}

\begin{center}
\begin{tabularx}{\textwidth}{@{}l Y@{}}
\toprule
\rowcolor{dermaSoft}
\textbf{Convención} & \textbf{Significado} \\
\midrule
\boton{Iniciar sesión} & Botón, enlace o casilla de la interfaz sobre el que se
hace clic. \\
\addlinespace
\campo{Correo electrónico} & Nombre de un campo de un formulario. \\
\addlinespace
\menu{Desarrolladores\flecha API} & Opción del menú de navegación; el símbolo
\flecha{} separa los niveles del menú. \\
\addlinespace
\mensaje{Credenciales incorrectas.} & Texto literal que muestra el sistema. \\
\addlinespace
\rol{Admin} & Rol de usuario. \\
\addlinespace
\texttt{/images} & Dirección (ruta) dentro de la plataforma. \\
\bottomrule
\end{tabularx}
\end{center}

Además, el manual utiliza los siguientes recuadros:

\begin{notabox}
Información complementaria útil para comprender el funcionamiento del sistema.
\end{notabox}
\begin{tipbox}
Recomendación práctica para trabajar de forma más cómoda o eficiente.
\end{tipbox}
\begin{warnbox}
Advertencia sobre acciones irreversibles o situaciones que requieren atención.
\end{warnbox}
\begin{adminbox}
Función disponible únicamente para cuentas con el rol Admin.
\end{adminbox}

% ==========================================================
\seccion{Descripción general del sistema}
\label{sec:descripcion}
% ==========================================================

\subsection{¿Qué es DermaUH Image Database?}

DermaUH Image Database es un repositorio web de imágenes dermatológicas clínicas
acompañadas de metadatos estandarizados: datos demográficos del paciente
(anonimizados), localización anatómica de la lesión, diagnóstico, método de
confirmación diagnóstica y, en los melanomas, datos histológicos. Su objetivo es
servir de base para la docencia, la investigación clínica y epidemiológica, y el
entrenamiento de modelos de Inteligencia Artificial en dermatología.

\subsection{Funcionalidades principales}

\begin{itemize}
  \item \textbf{Galería diagnóstica}: búsqueda y filtrado de imágenes por
  metadatos clínicos y consulta del detalle de cada caso.
  \item \textbf{Descarga de conjuntos de datos}: exportación de imágenes y/o
  metadatos en un archivo ZIP, sujeta a autorización.
  \item \textbf{Análisis estadístico}: panel con indicadores y gráficos sobre la
  composición del repositorio.
  \item \textbf{Instituciones}: consulta de los centros que aportan imágenes.
  \item \textbf{Gestión de cuentas}: registro, perfil personal y, para los
  administradores, aprobación de cuentas, autorización de descargas y
  gestión de roles.
  \item \textbf{Integración}: diccionario de metadatos y API documentada para
  aportar imágenes desde otras aplicaciones.
\end{itemize}

\subsection{Perfiles de usuario}
\label{sec:roles}

Las acciones disponibles dependen del \textbf{rol} de la cuenta. Toda cuenta nueva
recibe automáticamente el rol \rol{Viewer}; los administradores pueden asignar
roles adicionales, y una misma cuenta puede tener varios a la vez.

\begin{center}
\begin{tabularx}{\textwidth}{@{}l Y@{}}
\toprule
\rowcolor{dermaSoft}
\textbf{Rol} & \textbf{Descripción} \\
\midrule
\rol{Viewer} (Visualizador) & Explora la galería completa, consulta el detalle
de las imágenes, las instituciones y los análisis, gestiona su perfil y puede
solicitar autorización para descargar. \\
\addlinespace
\rol{Contributor} (Contribuidor) & Además de lo anterior, puede aportar imágenes al
repositorio mediante la API de integración (sección~\ref{sec:desarrolladores}). \\
\addlinespace
\rol{Admin} (Administrador) & Control total: aprueba cuentas y solicitudes de
descarga, gestiona usuarios y roles, elimina imágenes y descarga sin necesidad de
autorización previa. \\
\bottomrule
\end{tabularx}
\end{center}

La tabla siguiente resume los permisos de cada perfil, incluido el visitante que
navega sin haber iniciado sesión.

\begin{center}
\begin{tabularx}{\textwidth}{@{}Y c c c c@{}}
\toprule
\rowcolor{dermaSoft}
\textbf{Función} & \textbf{Visitante} & \textbf{Viewer} & \textbf{Contributor} & \textbf{Admin} \\
\midrule
Ver la página de inicio, el análisis estadístico y los términos & \checkmark & \checkmark & \checkmark & \checkmark \\
Explorar la galería & 10 primeras & \checkmark & \checkmark & \checkmark \\
Consultar el detalle de una imagen pública & \checkmark & \checkmark & \checkmark & \checkmark \\
Ver las instituciones colaboradoras & -- & \checkmark & \checkmark & \checkmark \\
Solicitar autorización de descarga & -- & \checkmark & \checkmark & -- \\
Descargar imágenes y metadatos & -- & Con autorización & Con autorización & \checkmark \\
Aportar imágenes vía API & -- & -- & \checkmark & \checkmark \\
Eliminar imágenes & -- & -- & -- & \checkmark \\
Aprobar cuentas y solicitudes de descarga & -- & -- & -- & \checkmark \\
Gestionar usuarios y roles & -- & -- & -- & \checkmark \\
\bottomrule
\end{tabularx}
\end{center}

\begin{notabox}
La galería muestra únicamente las imágenes marcadas como \textbf{públicas} en el
repositorio. Las imágenes no públicas no aparecen en los listados ni en las
descargas.
\end{notabox}

% ==========================================================
\seccion{Requisitos para el uso}
% ==========================================================

Para utilizar la plataforma solo se necesita:

\begin{itemize}
  \item Un \textbf{navegador web actualizado}: Google Chrome, Microsoft Edge,
  Mozilla Firefox o Safari en sus versiones recientes.
  \item \textbf{JavaScript habilitado} y el almacenamiento de datos del sitio
  permitido: la plataforma guarda la sesión en el navegador.
  \item Una \textbf{conexión a Internet estable}. La aplicación mantiene una
  conexión permanente con el servidor mientras la página está abierta; si se
  interrumpe, el navegador mostrará un aviso de reconexión.
  \item Una \textbf{cuenta de correo electrónico} propia, necesaria para
  registrarse, confirmar la cuenta y recuperar la contraseña.
  \item Una pantalla con resolución de al menos 1280\,×\,720 píxeles para
  visualizar con comodidad los gráficos. La interfaz se adapta también a tabletas
  y teléfonos.
\end{itemize}

No es necesario instalar ningún programa en el equipo del usuario.

% ==========================================================
\seccion{Primeros pasos}
\label{sec:primeros-pasos}
% ==========================================================

\subsection{Acceder a la plataforma}

Abra el navegador y escriba la dirección (URL) de la plataforma que le haya
facilitado el administrador de su institución. En una instalación local, la
dirección por defecto es \texttt{http://localhost:5130}. Se mostrará la página de
inicio (Figura~\ref{fig:inicio-anonimo}).

\captura{inicio-anonimo}{Página de inicio vista por un visitante sin sesión iniciada.}

\subsection{Elementos de la interfaz}

En la parte superior de todas las páginas se encuentra la \textbf{barra de
navegación}, con los siguientes accesos:

\begin{center}
\begin{tabularx}{\textwidth}{@{}l l Y@{}}
\toprule
\rowcolor{dermaSoft}
\textbf{Opción} & \textbf{Ruta} & \textbf{Contenido} \\
\midrule
\menu{Inicio} & \texttt{/} & Resumen de la plataforma y accesos rápidos. \\
\menu{Imágenes} & \texttt{/images} & Galería diagnóstica (sección~\ref{sec:galeria}). \\
\menu{Instituciones} & \texttt{/institutions} & Centros colaboradores (sección~\ref{sec:instituciones}). \\
\menu{Análisis} & \texttt{/analytics} & Panel estadístico (sección~\ref{sec:analisis}). \\
\menu{Desarrolladores\flecha Metadatos} & \texttt{/developers/metadata} & Diccionario de metadatos. \\
\menu{Desarrolladores\flecha API} & \texttt{/developers/api} & Documentación de la API. \\
\menu{Usuarios} & \texttt{/users} & Gestión de usuarios (solo \rol{Admin}). \\
\bottomrule
\end{tabularx}
\end{center}

A la derecha de la barra aparecen los botones \boton{Iniciar sesión} y
\boton{Registrarse} cuando no hay una sesión abierta. Tras iniciar sesión se
sustituyen por el \textbf{menú de usuario}: un botón con las iniciales y el nombre
de la persona que, al pulsarlo, despliega las opciones \boton{Mi perfil} y
\boton{Cerrar sesión} (Figura~\ref{fig:menu-usuario}). En pantallas estrechas el
menú se agrupa en un botón con tres líneas horizontales.

\captura[0.6\textwidth]{menu-usuario}{Menú de usuario desplegado tras iniciar sesión.}

\subsection{Crear una cuenta}
\label{sec:registro}

\begin{pasos}
  \item En la barra de navegación, pulse \boton{Registrarse}.
  \item Complete los campos \campo{Nombre}, \campo{Apellido} y
  \campo{Correo electrónico} (Figura~\ref{fig:registro}).
  \item Escriba una contraseña en \campo{Contraseña} y repítala en
  \campo{Confirmar contraseña}. Debe tener al menos 8 caracteres e incluir una
  letra mayúscula, una minúscula y un número.
  \item Lea los \hyperref[sec:terminos]{Términos y Condiciones} (el enlace se abre
  en una pestaña nueva) y marque la casilla de aceptación.
  \item Pulse \boton{Crear cuenta}. El sistema confirma el registro e indica que
  revise su correo.
\end{pasos}

\captura[0.8\textwidth]{registro}{Formulario de registro de una cuenta nueva.}

Una cuenta recién creada \textbf{todavía no permite iniciar sesión}. Para
activarla deben completarse dos condiciones, en cualquier orden:

\begin{enumerate}
  \item \textbf{Confirmar el correo electrónico} mediante el enlace recibido
  (sección~\ref{sec:confirmar}).
  \item \textbf{Aprobación de un administrador}, que recibe un aviso por correo
  cada vez que se registra una cuenta (sección~\ref{sec:aprobar-registros}).
  Cuando la cuenta es aprobada, usted recibe un correo de confirmación.
\end{enumerate}

\subsection{Confirmar el correo electrónico}
\label{sec:confirmar}

\begin{pasos}
  \item Abra el mensaje \mensaje{Confirma tu correo - DermaUH} recibido en su
  bandeja de entrada (revise también la carpeta de correo no deseado).
  \item Pulse el enlace de confirmación. Se abrirá la plataforma, que valida el
  enlace automáticamente.
  \item Si la validación es correcta se muestra la pantalla de la
  Figura~\ref{fig:confirmacion-correo}. Si el enlace no es válido o ha
  caducado, la página muestra un mensaje de error.
\end{pasos}

\captura{confirmacion-correo}{Confirmación de correo realizada correctamente.}

\subsection{Iniciar sesión}
\label{sec:login}

\begin{pasos}
  \item Pulse \boton{Iniciar sesión} en la barra de navegación.
  \item Introduzca su \campo{Correo electrónico} y su \campo{Contraseña}
  (Figura~\ref{fig:login}).
  \item Pulse el botón \boton{Iniciar sesión}. Si los datos son correctos, el
  menú de usuario aparece en la barra de navegación.
\end{pasos}

\captura{login}{Pantalla de inicio de sesión.}

\textbf{Acceso con Google.} Si la institución ha habilitado esta opción, debajo del
formulario aparece el botón de Google. Para usarlo, marque primero la casilla de
aceptación de los Términos y Condiciones y después pulse el botón. La primera vez
que se accede con Google se crea la cuenta y queda \textbf{pendiente de
aprobación} por un administrador; el correo de Google se considera ya confirmado.

\begin{notabox}
Por seguridad, la sesión caduca al cabo de un tiempo (por defecto, 24 horas). Al
caducar, el sistema le pedirá que vuelva a iniciar sesión.
\end{notabox}

\subsection{Cerrar sesión}

Pulse el menú de usuario y seleccione \boton{Cerrar sesión}. Es recomendable
hacerlo siempre al utilizar un equipo compartido.

\subsection{Recuperar la contraseña}
\label{sec:recuperar}

\begin{pasos}
  \item En la pantalla de inicio de sesión, pulse \boton{¿Olvidaste tu contraseña?}.
  \item Escriba su \campo{Correo electrónico} y pulse
  \boton{Enviar enlace} (Figura~\ref{fig:olvido-contrasena}).
  \item Abra el correo recibido y pulse el enlace para restablecer la contraseña.
  \item En la página que se abre (Figura~\ref{fig:restablecer-contrasena}),
  escriba la \campo{Nueva contraseña}, repítala en \campo{Confirmar contraseña}
  y pulse \boton{Restablecer contraseña}.
  \item Inicie sesión con la nueva contraseña.
\end{pasos}

\captura[0.8\textwidth]{olvido-contrasena}{Solicitud de recuperación de contraseña.}

\begin{notabox}
Por seguridad, el sistema muestra siempre el mismo mensaje tanto si el correo
existe como si no, y solo envía instrucciones a cuentas con el correo ya
confirmado.
\end{notabox}

\captura[0.8\textwidth]{restablecer-contrasena}{Formulario para establecer una nueva contraseña.}

% ==========================================================
\seccion{Página de inicio}
% ==========================================================

La página de inicio ofrece una vista general de la plataforma:

\begin{itemize}
  \item \textbf{Indicadores}: número de imágenes en el repositorio, instituciones
  activas y contribuidores activos.
  \item \textbf{Accesos directos}: \boton{Explorar galería} y
  \boton{Ver análisis estadístico}.
  \item \textbf{Tarjetas de funciones}: galería diagnóstica, instituciones y
  análisis estadístico. La última tarjeta depende del rol: los administradores
  ven \boton{Gestionar usuarios} y el resto de usuarios \boton{Ver mi perfil}
  (Figura~\ref{fig:inicio-tarjetas}).
\end{itemize}

\captura{inicio-tarjetas}{Tarjetas de acceso rápido de la página de inicio (vista de administrador).}

% ==========================================================
\seccion{Galería de imágenes}
\label{sec:galeria}
% ==========================================================

La galería diagnóstica (\menu{Imágenes}) permite explorar el repositorio y
seleccionar imágenes para su descarga. La pantalla se divide en el panel de
\textbf{filtros}, a la izquierda, y la zona de \textbf{resultados}, a la derecha
(Figura~\ref{fig:galeria}).

\captura{galeria}{Galería diagnóstica con el panel de filtros y los resultados.}

Cada tarjeta de resultado muestra la miniatura de la imagen, el nombre del archivo,
la categoría diagnóstica y el tipo de imagen. Encima de los resultados se indica
el número de casos encontrados.

\subsection{Consultar la galería sin iniciar sesión}

Un visitante sin sesión puede ver únicamente las \textbf{diez primeras} imágenes.
Al final de la lista se muestra el aviso \mensaje{Ha alcanzado el límite de vista
previa. Inicie sesión para explorar los N casos completos.} junto al botón
\boton{Iniciar sesión} (Figura~\ref{fig:galeria-anonimo}). Con la sesión iniciada,
las imágenes se cargan automáticamente a medida que se desplaza hacia abajo.

\captura[0.8\textwidth]{galeria-anonimo}{Límite de vista previa para visitantes sin sesión.}

\subsection{Filtrar las imágenes}

\begin{pasos}
  \item Para buscar por diagnóstico, escriba un término en el campo
  \campo{Diagnóstico} (por ejemplo, «melanoma»). Los resultados se actualizan
  mientras escribe.
  \item Marque una o varias casillas en los grupos de filtros. Dentro de un grupo,
  las opciones marcadas se combinan entre sí (se muestran los casos que cumplen
  cualquiera de ellas); los distintos grupos se aplican conjuntamente.
  \item Para eliminar todos los filtros, pulse \boton{Restablecer}.
\end{pasos}

\begin{center}
\begin{tabularx}{\textwidth}{@{}l Y@{}}
\toprule
\rowcolor{dermaSoft}
\textbf{Filtro} & \textbf{Opciones} \\
\midrule
Tipo de imagen & Dermoscópica; Clínica: Vista General; Clínica: Primer Plano; TBP: Vista
General; TBP: Primer Plano; RCM: Macroscópica; RCM: Tile; RCM: Mosaico. \\
Categoría & Benigno, Indeterminado, Maligno. \\
Tipo de lesión & Melanoma, Carcinoma Basocelular, Carcinoma Escamocelular, Otros. \\
Confirmación diagnóstica & Histopatología; Evaluación clínica individual; Imágenes
seriadas sin cambio; Consenso de expertos; Microscopía confocal. \\
Fototipo & I, II, III, IV, V, VI (escala de Fitzpatrick). \\
Sexo & Masculino, Femenino. \\
Sitio anatómico & Cabeza/Cuello; Extremidad Superior; Extremidad Inferior; Torso
Anterior; Torso Lateral; Torso Posterior; Palmas/Plantas; Oral/Genital. \\
\bottomrule
\end{tabularx}
\end{center}

\subsection{Seleccionar y descargar imágenes}
\label{sec:descargar}

Las descargas requieren haber iniciado sesión y, salvo para los administradores,
contar con autorización (sección~\ref{sec:autorizacion}).

\begin{pasos}
  \item Aplique los filtros que necesite.
  \item Marque la casilla situada en la esquina de cada imagen que desee
  descargar, o marque \boton{Seleccionar todo lo cargado} para seleccionar todas
  las imágenes visibles. El contador indica cuántas hay seleccionadas
  (Figura~\ref{fig:galeria-seleccion}).
  \item En el desplegable \campo{Formato}, elija el contenido del archivo:
  \textit{Imágenes y metadatos}, \textit{Solo metadatos} o \textit{Solo imágenes}.
  \item Pulse \boton{Descargar selección} para descargar las imágenes marcadas, o
  \boton{Descargar todo} para descargar \textbf{todos} los casos que cumplen los
  filtros actuales, aunque no estén cargados en pantalla.
  \item El navegador guarda un archivo ZIP.
\end{pasos}

\captura{galeria-seleccion}{Galería filtrada por categoría «Maligno» con tres imágenes seleccionadas.}

El archivo ZIP descargado contiene:

\begin{itemize}
  \item \texttt{metadata.csv}: una fila por imagen con todos sus metadatos
  clínicos (si el formato incluye metadatos). Puede abrirse con cualquier hoja de
  cálculo.
  \item Carpeta \texttt{images/}: los archivos de imagen, nombrados con su
  identificador público (si el formato incluye imágenes).
  \item \texttt{errores.txt}: solo aparece si alguna imagen no pudo incluirse, y
  enumera cuáles.
\end{itemize}

\subsection{Solicitar autorización de descarga}
\label{sec:autorizacion}

La primera vez que un usuario sin autorización intenta descargar, se abre la
ventana \textbf{Solicitar autorización para descargar}
(Figura~\ref{fig:modal-solicitud}).

\begin{pasos}
  \item En \campo{Motivo de uso}, describa para qué necesita las imágenes (por
  ejemplo, el proyecto de investigación o la asignatura).
  \item En \campo{Institución}, indique la institución a la que pertenece.
  \item Pulse \boton{Enviar solicitud}. Todos los administradores reciben un aviso
  por correo.
  \item Consulte el estado de la solicitud en \menu{Mi perfil}, apartado
  \textbf{Mis solicitudes de descarga} (sección~\ref{sec:mis-solicitudes}).
  Recibirá también un correo cuando sea aprobada o denegada.
\end{pasos}

\captura{modal-solicitud}{Ventana de solicitud de autorización de descarga.}

\begin{notabox}
Una vez aprobada una solicitud, la autorización queda asociada a la cuenta de
forma permanente: no es necesario solicitarla de nuevo para futuras descargas.
\end{notabox}

\subsection{Consultar el detalle de una imagen}
\label{sec:detalle}

Pulse sobre cualquier tarjeta de la galería para abrir su ficha de detalle
(Figura~\ref{fig:detalle-imagen}). La ficha es de solo lectura y muestra la imagen
ampliada junto a sus metadatos, agrupados en:

\begin{itemize}
  \item \textbf{Información general}: identificador público, archivo, tipo de
  imagen, manipulación, tipo dermatoscópico, contribuidor e institución.
  \item \textbf{Datos del paciente}: edad aproximada, sexo, color de piel,
  fototipo, antecedentes personales, antecedentes personales y familiares de
  melanoma, exposición solar y provincia.
  \item \textbf{Información clínica de la lesión}: sitio anatómico general y
  especial, y diámetro mayor en milímetros.
  \item \textbf{Diagnóstico}: diagnóstico clínico, diagnóstico histopatológico,
  categoría, tipo de lesión y método de confirmación.
  \item \textbf{Histología (melanoma)}: espesor de Breslow, índice mitótico y
  ulceración, cuando la lesión es un melanoma.
  \item \textbf{Notas}: notas clínicas y comentarios dermatoscópicos, si existen.
\end{itemize}

Para volver a la galería pulse \boton{Volver al listado}.

\captura[0.72\textwidth]{detalle-imagen}{Ficha de detalle de una imagen.}

% ==========================================================
\seccion{Instituciones}
\label{sec:instituciones}
% ==========================================================

La página \menu{Instituciones} presenta los centros que han aportado imágenes al
repositorio, en un carrusel de tarjetas que puede desplazarse lateralmente. Cada
tarjeta muestra el nombre, la descripción, el país y el número de imágenes
aportadas; en la parte superior se resumen el total de instituciones y de
imágenes (Figura~\ref{fig:instituciones}).

\captura{instituciones}{Instituciones colaboradoras.}

\begin{notabox}
El listado de instituciones solo se carga con la sesión iniciada. La página es de
consulta: las instituciones se registran automáticamente al aportar imágenes a
través de la API.
\end{notabox}

% ==========================================================
\seccion{Análisis estadístico}
\label{sec:analisis}
% ==========================================================

El panel \menu{Análisis} resume la composición del repositorio para apoyar la
investigación epidemiológica y la evaluación de la calidad de los datos. Es de
acceso público. Los gráficos se calculan con los datos disponibles en el momento
de la consulta; al pasar el puntero sobre un elemento se muestra su valor.

\subsection{Indicadores y distribución anatómica}

La parte superior muestra los indicadores clave (total de imágenes, imágenes
públicas, instituciones y contribuidores) y tres gráficos
(Figura~\ref{fig:analitica-resumen}):

\begin{itemize}
  \item \textbf{Mapeo de melanoma}: silueta corporal en la que cada región se
  colorea según el número de casos de melanoma localizados en ella.
  \item \textbf{Sitios anatómicos (general)}: número de imágenes por región
  anatómica.
  \item \textbf{Tendencia de cargas}: número de imágenes incorporadas cada mes.
\end{itemize}

\captura{analitica-resumen}{Indicadores clave, mapa corporal de melanoma y distribución anatómica.}

\subsection{Correlaciones clínicas}

Cuatro gráficos de barras agrupadas cruzan dos variables: Diagnóstico frente a
Sexo, Tipo de lesión frente a Grupo de edad, Diagnóstico frente a Color de piel y
Diagnóstico frente a Sitio anatómico (Figura~\ref{fig:analitica-correlaciones}).

\captura{analitica-correlaciones}{Correlaciones clínicas.}

\subsection{Distribución geográfica}

El mapa de Cuba colorea cada provincia según el número de imágenes procedentes de
ella, lo que permite valorar la representatividad territorial del repositorio
(Figura~\ref{fig:analitica-geografica}).

\captura{analitica-geografica}{Representatividad de las imágenes por provincia.}

\subsection{Distribuciones generales}

Ocho gráficos circulares muestran la proporción de casos según categoría
diagnóstica, tipo de lesión, sexo, color de piel, grupo de edad, fototipo, método
de confirmación y tipo de imagen (Figura~\ref{fig:analitica-distribuciones}).

\captura[0.85\textwidth]{analitica-distribuciones}{Distribuciones generales del repositorio.}

\subsection{Calidad de los datos}

El gráfico de completitud indica, para cada campo clínico, el porcentaje de
imágenes que tienen ese dato informado. Es útil para identificar qué metadatos
conviene completar (Figura~\ref{fig:analitica-completitud}).

\captura[0.85\textwidth]{analitica-completitud}{Completitud de los metadatos clínicos.}

% ==========================================================
\seccion{Mi perfil}
\label{sec:perfil}
% ==========================================================

Para abrir el perfil, pulse el menú de usuario y seleccione \boton{Mi perfil}. La
página requiere sesión iniciada (Figura~\ref{fig:perfil-usuario}).

A la izquierda se muestra el \textbf{resumen de la cuenta}: iniciales, nombre,
correo, estado de verificación del correo, roles asignados, fecha de alta y si la
cuenta está autorizada para descargas (por ejemplo, \mensaje{Sin autorización de
descarga masiva}). A la derecha se encuentran, de arriba abajo, el listado
\textbf{Mis solicitudes de descarga} y los formularios \textbf{Información
personal} y \textbf{Cambiar contraseña}, que se describen a continuación. Las
cuentas con rol Admin ven además, al principio, los apartados de aprobación
descritos en la sección~\ref{sec:admin}.

\captura[0.8\textwidth]{perfil-usuario}{Perfil de un usuario con rol Viewer.}

\subsection{Actualizar los datos personales}

\begin{pasos}
  \item En el apartado \textbf{Información personal}, modifique los campos
  \campo{Nombre}, \campo{Apellido} o \campo{Teléfono (opcional)}.
  \item Pulse \boton{Guardar cambios}.
\end{pasos}

El correo electrónico no puede modificarse, ya que identifica la cuenta.

\subsection{Cambiar la contraseña}

\begin{pasos}
  \item En el apartado \textbf{Cambiar contraseña}, escriba su
  \campo{Contraseña actual}.
  \item Escriba la \campo{Nueva contraseña} y repítala en
  \campo{Confirmar nueva contraseña}.
  \item Pulse \boton{Cambiar contraseña}.
\end{pasos}

\subsection{Seguimiento de las solicitudes de descarga}
\label{sec:mis-solicitudes}

El apartado \textbf{Mis solicitudes de descarga} lista las solicitudes enviadas,
con su fecha, motivo, institución y estado: \textit{Pendiente},
\textit{Aprobada} o \textit{Denegada}.

% ==========================================================
\seccion{Tareas de administración}
\label{sec:admin}
% ==========================================================

\begin{adminbox}
Las funciones de esta sección solo están disponibles para cuentas con el rol
Admin.
\end{adminbox}

\subsection{Aprobar o denegar registros}
\label{sec:aprobar-registros}

Las cuentas nuevas aparecen en \menu{Mi perfil}, apartado
\textbf{Solicitudes de registro pendientes} (Figura~\ref{fig:perfil-admin-registros}).

\begin{pasos}
  \item Abra \menu{Mi perfil}. Si hay muchas solicitudes, escriba parte del correo
  en \campo{Buscar por email} y pulse \boton{Buscar}.
  \item Compruebe el nombre y el correo del solicitante.
  \item Pulse \boton{Aceptar} para activar la cuenta, o \boton{Denegar} para
  rechazarla. En ambos casos el solicitante recibe un correo.
\end{pasos}

\captura{perfil-admin-registros}{Solicitudes de registro pendientes.}

\begin{warnbox}
Denegar una solicitud \textbf{elimina la cuenta}. Si la persona desea acceder más
adelante, deberá registrarse de nuevo.
\end{warnbox}

\begin{notabox}
Aprobar una cuenta no confirma su correo: el usuario solo podrá iniciar sesión
cuando, además, haya pulsado el enlace de confirmación que recibió al
registrarse.
\end{notabox}

\subsection{Revisar solicitudes de descarga}

Las solicitudes de autorización de descarga aparecen en \menu{Mi perfil},
apartado \textbf{Solicitudes de descarga pendientes}, con el solicitante, su
correo, el motivo, la institución y la fecha (Figura~\ref{fig:perfil-admin-descargas}).

\begin{pasos}
  \item Lea el motivo de uso y la institución indicados.
  \item Pulse \boton{Aceptar} para conceder la autorización de forma permanente, o
  \boton{Denegar} para rechazarla. El usuario recibe un correo con la decisión.
\end{pasos}

\captura{perfil-admin-descargas}{Solicitudes de descarga pendientes de revisión.}

\subsection{Gestionar usuarios}

La página \menu{Usuarios} lista todas las cuentas con su nombre, correo, roles,
estado y fecha de registro (Figura~\ref{fig:usuarios}). El listado puede
filtrarse con \campo{Buscar por email}, \campo{Rol} (Todos, Admin, Contributor,
Viewer o Sin rol) y \campo{Estado} (Todos, Activos o Inactivos).

\captura{usuarios}{Gestión de usuarios.}

\subsubsection*{Crear un usuario}

\begin{pasos}
  \item Pulse \boton{+ Nuevo usuario}. Se despliega el formulario de la
  Figura~\ref{fig:usuarios-nuevo}.
  \item Complete \campo{Nombre}, \campo{Apellido}, \campo{Email} y
  \campo{Contraseña} (todos obligatorios).
  \item Pulse \boton{Guardar}. Para cerrar el formulario sin crear la cuenta, pulse
  \boton{Cancelar}.
\end{pasos}

\captura{usuarios-nuevo}{Formulario de creación de un usuario.}

La cuenta se crea con el rol \rol{Viewer} y sigue el mismo circuito que un registro
normal: la persona recibe el correo de confirmación y la cuenta aparece en las
solicitudes de registro pendientes hasta que se aprueba o se activa.

\subsubsection*{Asignar o quitar roles}

\begin{pasos}
  \item Localice el usuario en el listado.
  \item Para añadir un rol, elíjalo en el primer selector de la columna
  \textbf{Acciones} (solo se ofrecen los roles que aún no tiene) y pulse
  \boton{Asignar}.
  \item Para quitar un rol, elíjalo en el segundo selector (roles asignados) y pulse
  \boton{Quitar}.
\end{pasos}

Los cambios de rol se aplican la próxima vez que el usuario inicie sesión.

\subsubsection*{Activar, desactivar o eliminar una cuenta}

\begin{itemize}
  \item \boton{Desactivar} impide que la cuenta inicie sesión sin borrarla;
  \boton{Activar} la habilita de nuevo. El navegador pide confirmación.
  \item \boton{Eliminar} borra la cuenta tras confirmar el mensaje
  \mensaje{¿Deseas eliminar este usuario? Esta acción no se puede deshacer.}
\end{itemize}

\subsection{Eliminar una imagen}

\begin{pasos}
  \item Abra la ficha de detalle de la imagen (sección~\ref{sec:detalle}).
  \item Pulse \boton{Eliminar} en la parte superior derecha
  (Figura~\ref{fig:detalle-admin}).
  \item Confirme el mensaje \mensaje{¿Deseas eliminar esta imagen? Esta acción no se
  puede deshacer.} El sistema vuelve a la galería.
\end{pasos}

\captura{detalle-admin}{Ficha de detalle vista por un administrador, con el botón «Eliminar».}

\begin{warnbox}
La eliminación de imágenes y de usuarios no se puede deshacer desde la interfaz.
\end{warnbox}

% ==========================================================
\seccion{Sección para desarrolladores}
\label{sec:desarrolladores}
% ==========================================================

El menú \menu{Desarrolladores} reúne la información necesaria para integrar otras
aplicaciones con la plataforma. Es de acceso público.

\subsection{Diccionario de metadatos}

\menu{Desarrolladores\flecha Metadatos} describe cada campo clínico del
repositorio: nombre, significado, tipo de dato y valores admitidos
(Figura~\ref{fig:desarrolladores-metadatos}). El botón
\boton{Descargar diccionario (CSV)} exporta la lista completa. El mismo
diccionario está disponible en la ruta \texttt{/metadata}.

\captura{desarrolladores-metadatos}{Diccionario de metadatos.}

\subsection{API de carga de imágenes}

\menu{Desarrolladores\flecha API} documenta cómo aportar imágenes desde otra
aplicación (Figura~\ref{fig:desarrolladores-api}). Los datos esenciales son:

\begin{center}
\begin{tabularx}{\textwidth}{@{}l Y@{}}
\toprule
\rowcolor{dermaSoft}
\textbf{Elemento} & \textbf{Detalle} \\
\midrule
Operación & \texttt{POST /api/images} \\
Autenticación & Token \textit{Bearer} (JWT) obtenido con \texttt{POST /api/auth/login} \\
Roles permitidos & \rol{Admin} o \rol{Contributor} \\
Formato & \texttt{multipart/form-data} (archivo de imagen y campos de metadatos) \\
Campos obligatorios & Archivo, identificador del contribuidor, edad aproximada,
sexo, sitio anatómico general, diagnóstico clínico y exposición solar \\
\bottomrule
\end{tabularx}
\end{center}

La página detalla cada campo con su tipo, obligatoriedad y reglas de validación;
por ejemplo, el tipo de lesión solo se admite si la categoría diagnóstica es
\textit{Maligno}, y los datos histológicos solo si la lesión es un melanoma.

\captura[0.8\textwidth]{desarrolladores-api}{Documentación de la API de carga de imágenes.}

% ==========================================================
\seccion{Términos y condiciones}
\label{sec:terminos}
% ==========================================================

La página de Términos y Condiciones (\texttt{/terms}) explica el origen de las
imágenes, la existencia de consentimiento informado de los pacientes, los usos
permitidos del contenido (docencia, investigación y mejora de la atención
sanitaria) y las garantías de anonimización de los datos. Su aceptación es
obligatoria para registrarse y para acceder con Google. Es recomendable leerla
antes de solicitar una autorización de descarga (Figura~\ref{fig:terminos}).

\captura[0.8\textwidth]{terminos}{Página de Términos y Condiciones de Uso.}

% ==========================================================
\seccion{Mensajes del sistema y solución de problemas}
\label{sec:problemas}
% ==========================================================

\subsection{Mensajes del sistema}

\begin{longtable}{@{}>{\raggedright\arraybackslash}p{0.34\textwidth} >{\raggedright\arraybackslash}p{0.6\textwidth}@{}}
\toprule
\rowcolor{dermaSoft}
\textbf{Mensaje} & \textbf{Causa y solución} \\
\midrule
\endhead
\mensaje{Credenciales incorrectas.} & El correo o la contraseña no son correctos.
Revíselos (las mayúsculas cuentan) o recupere la contraseña
(sección~\ref{sec:recuperar}). \\
\addlinespace
\mensaje{Debes confirmar tu correo electrónico antes de iniciar sesión.} & No se
ha pulsado el enlace del correo de confirmación. Búsquelo en su bandeja de
entrada o en correo no deseado. \\
\addlinespace
\mensaje{Tu cuenta no está activa. Espere a que sea aprobada o contacte al
administrador.} & La cuenta aún no ha sido aprobada, o un administrador la ha
desactivado. \\
\addlinespace
\mensaje{Tu cuenta ha sido creada y está pendiente de aprobación por un
administrador.} & Primer acceso con Google. Espere el correo de aprobación. \\
\addlinespace
\mensaje{Ya existe una cuenta con ese correo electrónico.} & El correo ya está
registrado. Inicie sesión o recupere la contraseña. \\
\addlinespace
\mensaje{Debe aceptar los términos y condiciones antes de iniciar sesión con
Google.} & Marque la casilla de aceptación antes de pulsar el botón de Google. \\
\addlinespace
\mensaje{Ha alcanzado el límite de vista previa\ldots} & Los visitantes solo ven
diez imágenes. Inicie sesión para ver el catálogo completo. \\
\addlinespace
\mensaje{No se encontraron imágenes con los filtros actuales.} & Ningún caso
cumple todos los filtros. Pulse \boton{Restablecer} o reduzca los filtros. \\
\addlinespace
\mensaje{Necesita autorización de un administrador para descargar contenido.} &
Aparece en la ventana de solicitud: su cuenta aún no tiene autorización de
descarga. Envíe la solicitud (sección~\ref{sec:autorizacion}). \\
\addlinespace
\mensaje{No fue posible sincronizar las métricas en este momento.} & El servidor
no respondió al cargar los indicadores. Recargue la página pasados unos minutos. \\
\bottomrule
\end{longtable}

\subsection{Problemas frecuentes}

\textbf{No recibo el correo de confirmación o de recuperación.}\\
Revise la carpeta de correo no deseado y compruebe que escribió bien la
dirección. Si el problema persiste, contacte con el administrador: es posible que
el envío de correo no esté configurado en la instalación.

\textbf{La página muestra un aviso de reconexión o deja de responder.}\\
La conexión con el servidor se ha interrumpido. Espere a que se restablezca o
recargue la página con la tecla F5.

\textbf{El enlace del correo indica que no es válido.}\\
Los enlaces de confirmación y de recuperación caducan (por defecto, a las 24
horas). Para la contraseña, solicite un enlace nuevo; para la confirmación,
contacte con el administrador.

\textbf{La página de instituciones aparece vacía.}\\
El listado de instituciones requiere iniciar sesión. Sin sesión se muestra
\mensaje{No se encontraron instituciones con imágenes aportadas.}

\textbf{Al recargar \menu{Mi perfil} o \menu{Usuarios} aparece la pantalla de inicio
de sesión.}\\
Ocurre al recargar la página (F5) o al abrir su dirección directamente, aunque la
sesión siga abierta (su nombre continúa en la barra de navegación). Acceda a estas
páginas desde el menú de usuario o desde la barra de navegación.

\textbf{El mapa de provincias muestra marcas de agua sobre el fondo.}\\
El proveedor de mapas base requiere una clave de acceso que aún no está
configurada. Los datos por provincia (marcadores y tabla lateral) se muestran
correctamente.

\textbf{Me han asignado un rol nuevo pero no veo las nuevas opciones.}\\
Cierre sesión y vuelva a iniciarla para que se apliquen los permisos.

% ==========================================================
\seccion{Preguntas frecuentes}
% ==========================================================

\textbf{¿Por qué no puedo iniciar sesión justo después de registrarme?}\\
Porque la cuenta necesita que confirme su correo y que un administrador la
apruebe (sección~\ref{sec:registro}).

\textbf{¿Cómo subo nuevas imágenes al repositorio?}\\
La interfaz web no incluye un formulario de carga. Las imágenes se aportan mediante
la API (sección~\ref{sec:desarrolladores}), con una cuenta de rol Contributor o
Admin, normalmente desde otra aplicación del ecosistema DermaUH.

\textbf{¿Puedo modificar los metadatos de una imagen?}\\
No desde la interfaz web. La corrección de metadatos está reservada a los
administradores a través de la API.

\textbf{¿Tengo que pedir autorización cada vez que descargo?}\\
No. Una vez aprobada, la autorización es permanente para su cuenta.

\textbf{¿Qué contiene exactamente el archivo descargado?}\\
Un ZIP con \texttt{metadata.csv} y/o la carpeta \texttt{images/}, según el formato
elegido (sección~\ref{sec:descargar}).

\textbf{¿Por qué algunas imágenes del repositorio no aparecen en la galería?}\\
La galería solo muestra las imágenes marcadas como públicas.

% ==========================================================
\seccion{Glosario}
\label{sec:glosario}
% ==========================================================

\begin{description}[style=nextline, leftmargin=1.5em, itemsep=1pt, topsep=2pt]
  \item[Breslow (espesor de)] Profundidad de invasión de un melanoma, medida en
  milímetros en el estudio histológico. Es un factor pronóstico fundamental.
  \item[Carcinoma basocelular] Tipo de cáncer de piel más frecuente, de
  crecimiento lento y rara vez metastásico.
  \item[Categoría diagnóstica] Clasificación de la lesión como benigna,
  indeterminada o maligna.
  \item[Consentimiento informado] Autorización documentada del paciente para el
  uso de sus imágenes con fines asistenciales, docentes o de investigación.
  \item[CSV] Formato de texto con valores separados por comas que puede abrirse con
  hojas de cálculo.
  \item[Dermatoscopia] Técnica de exploración de la piel con un instrumento de
  aumento e iluminación que permite ver estructuras no visibles a simple vista.
  Puede ser con luz polarizada o no polarizada, y con o sin contacto.
  \item[Fototipo] Clasificación de Fitzpatrick (I a VI) de la respuesta de la piel a
  la exposición solar.
  \item[Histopatología] Estudio microscópico de una muestra de tejido; es el
  método de referencia para confirmar un diagnóstico.
  \item[Índice mitótico] Número de células en división por milímetro cuadrado en
  un melanoma.
  \item[JWT] \textit{JSON Web Token}: credencial temporal que la plataforma emite al
  iniciar sesión y que autoriza las peticiones a la API.
  \item[Melanoma] Cáncer de piel originado en los melanocitos; el de mayor
  gravedad.
  \item[Metadatos] Datos que describen cada imagen: paciente, lesión, diagnóstico
  y características de la captura.
  \item[RCM] Microscopía confocal de reflectancia: técnica de imagen no invasiva de
  alta resolución.
  \item[Rol] Perfil de permisos de una cuenta (Viewer, Contributor, Admin).
  \item[TBP] Fotografía corporal total (\textit{Total Body Photography}).
  \item[ZIP] Formato de archivo comprimido que agrupa varios archivos en uno.
\end{description}

% ==========================================================
\seccion{Soporte}
% ==========================================================

Para cualquier incidencia no resuelta en este manual (cuentas bloqueadas,
solicitudes pendientes durante mucho tiempo, errores de la aplicación), contacte
con el administrador de la plataforma de su institución. Al comunicar un problema,
indique:

\begin{itemize}
  \item el correo de su cuenta;
  \item la página en la que se encontraba y la acción que intentaba realizar;
  \item el texto exacto del mensaje mostrado, o una captura de pantalla;
  \item la fecha y hora aproximadas.
\end{itemize}

% ==========================================================
\clearpage
\appendix
\renewcommand{\thesection}{\Alph{section}}
\section{Anexo: instalación y despliegue con Docker}
\label{anexo:despliegue}
% ==========================================================

Este anexo está dirigido al personal técnico responsable de instalar la
plataforma. Los archivos mencionados se encuentran en la raíz del repositorio del
proyecto.

\subsection{Arquitectura}

La plataforma se despliega con \textbf{Docker Compose} como un conjunto de
contenedores:

\begin{center}
\begin{tabularx}{\textwidth}{@{}l l Y@{}}
\toprule
\rowcolor{dermaSoft}
\textbf{Servicio} & \textbf{Puerto (host)} & \textbf{Función} \\
\midrule
\texttt{web} & 5130 & Aplicación web (Blazor Server) que usan los usuarios. \\
\texttt{api} & 5131 & API REST. La usa la aplicación web y las integraciones externas. \\
\texttt{postgres} & 5432 (solo local) & Base de datos PostgreSQL 16. \\
\texttt{mailpit} & 8025 & Opcional: servidor de correo de pruebas. \\
\bottomrule
\end{tabularx}
\end{center}

Los datos persistentes se guardan en volúmenes de Docker:
\texttt{postgres\_data} (base de datos), \texttt{api\_images} (archivos de
imagen) y \texttt{api\_keys} / \texttt{web\_keys} (claves de protección de
datos, de las que dependen los enlaces de confirmación y de recuperación
enviados por correo).

\subsection{Requisitos}

\begin{itemize}
  \item Docker Engine con el complemento Docker Compose v2.
  \item Al menos 2\,GB de memoria y espacio en disco suficiente para las imágenes.
  \item Acceso a Internet durante la construcción de las imágenes (descarga de
  dependencias).
\end{itemize}

\subsection{Configuración}

\begin{pasos}
  \item Copie el archivo de ejemplo de configuración:
\begin{consola}
cp .env.example .env
\end{consola}
  \item Edite \texttt{.env} y ajuste, como mínimo, las variables de la tabla
  siguiente.
\end{pasos}

\begin{longtable}{@{}>{\raggedright\arraybackslash}p{0.36\textwidth} >{\raggedright\arraybackslash}p{0.58\textwidth}@{}}
\toprule
\rowcolor{dermaSoft}
\textbf{Variable} & \textbf{Descripción} \\
\midrule
\endhead
\texttt{POSTGRES\_PASSWORD} & Contraseña de la base de datos. \\
\texttt{JWT\_SECRET\_KEY} & Clave de firma de las sesiones; cadena aleatoria de al
menos 32 caracteres (\texttt{openssl rand -base64 48}). \\
\texttt{BOOTSTRAP\_ADMIN\_EMAIL}, \texttt{BOOTSTRAP\_ADMIN\_PASSWORD} & Cuenta del
primer administrador, que se crea al arrancar la API. \\
\texttt{FRONTEND\_BASE\_URL} & URL pública de la aplicación web; se usa en los
enlaces enviados por correo. \\
\texttt{SMTP\_*} & Servidor de correo saliente. Sin él, los usuarios no reciben
los correos de confirmación. \\
\texttt{GOOGLE\_CLIENT\_ID} & Opcional: habilita el acceso con Google. \\
\texttt{WEB\_PORT}, \texttt{API\_PORT} & Puertos publicados en el servidor. \\
\bottomrule
\end{longtable}

\begin{warnbox}
Antes de un despliegue en producción, sustituya los valores de ejemplo de
\texttt{POSTGRES\_PASSWORD}, \texttt{JWT\_SECRET\_KEY} y
\texttt{BOOTSTRAP\_ADMIN\_PASSWORD} por valores propios y seguros. El archivo
\texttt{.env} no debe incluirse en el control de versiones.
\end{warnbox}

\subsection{Puesta en marcha}

\begin{pasos}
  \item Construya las imágenes y arranque los servicios:
\begin{consola}
docker compose up -d --build
\end{consola}
  \item Compruebe que la API ha arrancado y aplicado las migraciones:
\begin{consola}
docker compose logs -f api
\end{consola}
  \item Abra \texttt{http://localhost:5130} (o la URL pública configurada) e
  inicie sesión con la cuenta del administrador inicial.
\end{pasos}

La API aplica automáticamente las migraciones de la base de datos en cada arranque;
no se requiere ningún paso manual.

\subsection{Primer administrador}

Una base de datos nueva no contiene usuarios. Si se definen
\texttt{BOOTSTRAP\_ADMIN\_EMAIL} y \texttt{BOOTSTRAP\_ADMIN\_PASSWORD}, la API crea
esa cuenta (activa y con el correo confirmado) y le asigna el rol Admin; si la
cuenta ya existe, solo se asegura de que tenga dicho rol. Una vez creado el
administrador pueden dejarse vacías.

\subsection{Correo electrónico}

Configure las variables \texttt{SMTP\_*} con los datos de su servidor de correo.
Para pruebas locales puede usarse Mailpit, que captura todos los correos sin
enviarlos:

\begin{consola}
# En .env: SMTP_SERVER_ADDRESS=mailpit  y  SMTP_SERVER_PORT=1025
docker compose --profile mail up -d
\end{consola}

Los correos capturados se consultan en \texttt{http://localhost:8025}.

\subsection{Operación y mantenimiento}

\begin{center}
\begin{tabularx}{\textwidth}{@{}l Y@{}}
\toprule
\rowcolor{dermaSoft}
\textbf{Tarea} & \textbf{Comando} \\
\midrule
Detener & \texttt{docker compose down} \\
Iniciar & \texttt{docker compose up -d} \\
Actualizar versión & \texttt{git pull} y \texttt{docker compose up -d --build} \\
Ver registros & \texttt{docker compose logs -f api web} \\
Estado de servicios & \texttt{docker compose ps} \\
\bottomrule
\end{tabularx}
\end{center}

\textbf{Copias de seguridad.} Deben respaldarse la base de datos y los archivos de
imagen:

\begin{consola}
docker compose exec postgres pg_dump -U root DermaImage > dermauh.sql
docker run --rm -v dermauh-image-database_api_images:/data \
  -v "$PWD":/backup alpine tar czf /backup/imagenes.tar.gz -C /data .
\end{consola}

\begin{warnbox}
\texttt{docker compose down -v} elimina también los volúmenes, es decir, la base
de datos y las imágenes. No lo utilice sin disponer de una copia de seguridad.
\end{warnbox}

\subsection{Publicación en Internet}

Para exponer la plataforma públicamente se recomienda situar delante de los
servicios \texttt{web} y \texttt{api} un proxy inverso (por ejemplo, Nginx o
Traefik) que proporcione HTTPS. El proxy debe \textbf{permitir conexiones
WebSocket} hacia \texttt{web}, ya que la aplicación mantiene una conexión
permanente con el navegador. Actualice \texttt{FRONTEND\_BASE\_URL} con la URL
pública definitiva.

\end{document}
